%% glimm_uhf_lean4_2026-07-23.tex
%%
%% A machine-checked construction of the UHF algebra M_{3^infty} and its
%% unique faithful trace, with a formalized proof of Glimm simplicity, in
%% Lean 4 over mathlib.
%%
%% Every mathematical claim below is either a kernel-verified Lean 4 theorem
%% at HEAD 91cc1a01 (2026-07-23), axioms [propext, Classical.choice,
%% Quot.sound], zero project axioms, zero sorries, or standard cited
%% literature. Verified build data: lake build PF = 4472 jobs, lake build
%% (full) = 8930 jobs, mathlib pin eed770a4, toolchain
%% leanprover/lean4:v4.24.0-rc1.

\documentclass[11pt,reqno]{amsart}

\usepackage[T1]{fontenc}
\usepackage{lmodern}
\usepackage{amsmath,amssymb,amsthm}
\usepackage{mathtools}
\usepackage{microtype}
\usepackage{geometry}
\geometry{margin=1in}
\usepackage[hidelinks]{hyperref}
\usepackage{enumitem}
\usepackage{booktabs}
\usepackage{seqsplit}

\emergencystretch=3em
\hbadness=6000
\sloppy

\theoremstyle{plain}
\newtheorem{theorem}{Theorem}[section]
\newtheorem{lemma}[theorem]{Lemma}
\newtheorem{proposition}[theorem]{Proposition}
\newtheorem{corollary}[theorem]{Corollary}

\theoremstyle{definition}
\newtheorem{definition}[theorem]{Definition}

\theoremstyle{remark}
\newtheorem{remark}[theorem]{Remark}

\newcommand{\lt}[1]{\texttt{\seqsplit{#1}}}
\newcommand{\Tinf}{T_\infty}
\newcommand{\tuhf}{\tau_{\mathrm{UHF}}}
\newcommand{\op}{\mathrm{op}}
\DeclareMathOperator{\Tr}{tr}
\DeclareMathOperator{\Trace}{trace}
\DeclareMathOperator{\diag}{diag}
\let\Re\relax
\DeclareMathOperator{\Re}{Re}

\title[A machine-checked UHF algebra and its unique faithful trace]{A
machine-checked construction of the UHF algebra
\texorpdfstring{$M_{3^\infty}$}{M(3\textasciicircum\ensuremath{\infty})}
and its unique faithful trace}

\author{Pablo Cohen}
\address{Mesa, Arizona, USA}
\email{psolorzano@gmail.com}
\urladdr{ORCID:~\href{https://orcid.org/0009-0002-0734-5565}{0009-0002-0734-5565}}

\date{July 23, 2026}

\begin{document}

\begin{abstract}
We report a Lean~4 formalization, over mathlib, of the uniformly
hyperfinite $C^*$-algebra $M_{3^\infty}$ and its canonical trace. The
algebra is built as the metric completion $\Tinf$ of the inductive limit
of the matrix tower $M_{3^k}(\mathbb{C})$ under the unital $*$-embeddings
$A \mapsto A \otimes I_3$, and carries mathlib's \texttt{CStarAlgebra}
instance. On it we construct a tracial state $\tuhf$ and prove, all
kernel-verified, that it is faithful: $\tuhf(x^*x) = 0$ implies $x = 0$ for
every $x \in \Tinf$. As a corollary $\Tinf$ is simple, every two-sided
ideal being $\{0\}$ or the whole algebra, which is Glimm's 1960 theorem for
this algebra. We further prove that $\tuhf$ is the unique continuous unital
tracial functional on $\Tinf$. Faithfulness, simplicity, and uniqueness of
the trace together identify $\Tinf$ as the Glimm $3^\infty$ factor. To the
best of our knowledge this is the first machine-checked construction of an
infinite-dimensional simple $C^*$-algebra with a canonical faithful and
unique trace. Every headline theorem depends only on the three standard
kernel axioms $[\texttt{propext}, \texttt{Classical.choice},
\texttt{Quot.sound}]$, with no project-level axioms and no \texttt{sorry},
verified at pinned toolchain and mathlib revisions.
\end{abstract}

\maketitle

%% =====================================================================
\section{Introduction}\label{sec:intro}
%% =====================================================================

The uniformly hyperfinite (UHF) $C^*$-algebras are among the oldest and
best-understood objects in operator algebra theory. They are the
norm-closures of increasing unions of full matrix algebras. Glimm
classified them in 1960 by a supernatural number and proved that each one
is simple and carries a unique tracial state, necessarily
faithful~\cite{glimm1960}. The standard references treat all of this as
foundational~\cite{davidson1996,blackadar2006,takesaki1979}. The algebra
of the present paper is $M_{3^\infty}$: the tower $M_{3^k}(\mathbb{C})$ with
the embeddings $A \mapsto A \otimes I_3$, completed in norm.

As a formalization target the object is attractive precisely because it is
concrete. It is a single, explicitly presented $C^*$-algebra whose three
defining structural facts, faithfulness of the canonical trace, simplicity,
and uniqueness of the trace, are theorems rather than definitions, and each
of the three requires a genuine analytic or algebraic argument to
establish. A proof assistant either has these theorems or it does not.
Before this work no proof assistant did. Formalized operator algebra is
young; mathlib, the Lean~4 mathematical library~\cite{mathlib2020,
moura-ullrich2021}, provides a $C^*$-algebra typeclass hierarchy, an
$L^2$-operator-norm structure on matrices, ring direct limits, metric
completions, and a continuous functional calculus, following the
functional-analysis program of Dupuis, Lewis, and
Macbeth~\cite{dupuis-lewis-macbeth2022}. It contains no concrete
infinite-dimensional simple $C^*$-algebra, no tracial state theory, no
conditional expectations, and no completely positive maps. Building
$M_{3^\infty}$ and proving the classical facts therefore means assembling,
from mathlib's primitives, a body of infrastructure that the pen-and-paper
literature takes for granted, and at several points finding proofs
different from the classical ones.

\subsection*{Why faithfulness is the hard property}

Everything about $\tuhf$ except faithfulness transfers from the matrix
level to the completion by soft means. Additivity, traciality, Hermitian
symmetry, and positivity are all closed conditions: each says that a
continuous function vanishes, or is nonnegative, on $\Tinf$, so it is
enough to check it on the dense image of the direct limit and pass to the
closure. Faithfulness is different in kind. The statement
$\tuhf(x^*x) = 0 \Rightarrow x = 0$ is an implication with a nontrivial
conclusion, and the set of $x$ satisfying it is not the preimage of a
closed set. A density argument cannot reach it: a limit of matrices, each
detected by its own trace, could in principle converge to a nonzero element
whose trace-of-square vanishes. Ruling this out needs global structure.
Classically one has two routes. Either faithfulness comes from simplicity,
because the left kernel of a state generates a closed two-sided ideal that
must be trivial when none exists, with simplicity supplied by Glimm's
theorem; or it comes from the norm-one conditional expectations onto the
finite levels, in the tradition of Tomiyama~\cite{tomiyama1957}, whose
contractivity is classically a consequence of complete positivity and the
Kadison--Schwarz inequality~\cite{kadison1952}. Neither route was available
in mathlib.

\subsection*{How the proof was actually found}

The development did not proceed cleanly, and the record is worth stating.
An early version of this material carried a completion-level
``faithfulness'' declaration that, on inspection, was a placeholder:
a named proposition equal to \texttt{True}, discharged by
\texttt{trivial}, together with a theorem whose hypothesis was the
conclusion restated. Neither proves anything, and a companion note demoted
completion-level faithfulness from conditional to open before the argument
was rebuilt.

The first serious attempt reduced faithfulness, in kernel-verified form, to
a single hypothesis: that every nonzero two-sided ideal of $\Tinf$ contains
an element of nonzero trace. This looked like progress. A four-fold
equivalence then showed that the hypothesis is not weaker than the goal but
logically equivalent to it, and to simplicity as well. The reduction had
restated the summit, not descended from it. The equivalence was worth
formalizing anyway, because it is an exact statement of where the classical
hard input has to enter, and it ruled out any attempt to reach the result
from something strictly softer.

What eventually closed the argument was the conditional-expectation route,
built from scratch with elementary replacements for the missing theory.
The trace $2$-norm contraction of the level-$k$ expectation, classically a
corollary of Kadison--Schwarz for a completely positive map, reduced to a
three-line finite Cauchy--Schwarz once the target inequality was seen to be
scalar rather than operator-valued. The operator-norm contraction,
classically the statement that a unital completely positive map has norm
one, followed from the isometry decomposition
$\mathcal{E}B = \tfrac{1}{3}\sum_t W_t^* B\, W_t$ together with the
$C^*$-identity and submultiplicativity, with no complete-positivity theory
at all. With both contractions in hand, faithfulness was proved directly,
and simplicity was obtained as a corollary through the equivalence above.
This inverts the classical order, in which simplicity is proved first and
faithfulness deduced from it.

\subsection*{A note on names}

The formalization grew out of a larger exploratory project, and the results
presented here are self-contained and depend on none of that project's
conjectural content. Lean identifiers retain their original names: the
direct limit is \lt{TimelessFieldRing} and its completion
\lt{TimelessFieldCompletion} in the source, and many declarations carry a
\texttt{substrate} or version prefix. These are names, nothing more. In this
paper the objects are written $\varinjlim_k A_k$ and $\Tinf$, and every
property claimed of them is one of the kernel-verified theorems named below
or a standard cited fact.

\subsection*{Roadmap}

Section~\ref{sec:tower} constructs the tower, its inductive limit, and the
completion, and registers the $C^*$-algebra instance.
Section~\ref{sec:trace} builds $\tuhf$ and its soft properties.
Section~\ref{sec:engine} develops the algebraic engine: a Cauchy--Schwarz
inequality for the trace pairing, the trace-null ideal, the Weyl averaging
computation, simplicity below the completion, and the four-fold
equivalence. Section~\ref{sec:condexp} is the analytic core: the level
expectation, its two contraction estimates, the extension to the
completion, and the proof of faithfulness, followed by simplicity,
uniqueness of the trace, and the identification of $\Tinf$ as a factor.
Section~\ref{sec:repro} records the toolchain and verification data, and
Section~\ref{sec:related} discusses related work.

%% =====================================================================
\section{The tower, the limit, and the completion}\label{sec:tower}
%% =====================================================================

At level $k \in \mathbb{N}$ set
\[
A_k \;:=\; \mathrm{Matrix}\,(\mathrm{Fin}\,3^k)\,(\mathrm{Fin}\,3^k)\,\mathbb{C},
\]
a finite-dimensional $C^*$-algebra under the $L^2$-operator norm. This is
mathlib's scoped instance \lt{Matrix.Norms.L2Operator}, which supplies the
$C^*$-identity $\|M^*M\| = \|M\|^2$. Successive levels embed by a ternary
tensor construction,
\[
\iota_k \colon A_k \hookrightarrow A_{k+1}, \qquad
A \longmapsto \mathrm{reindex}\,(\lt{levelStepEquiv}\,k)\,(A \otimes I_3),
\]
where \lt{levelStepEquiv}~$k$ is the equivalence
$\mathrm{Fin}\,3^k \times \mathrm{Fin}\,3 \simeq \mathrm{Fin}\,3^{k+1}$
coming from the ternary decomposition of $3^{k+1}$, and $\otimes$ is the
Kronecker product. In the source this map is \lt{substrateEmbedMatrix},
with iterated form \lt{substrateRingHomIter}. Each $\iota_k$ is a unital
$*$-ring homomorphism and an isometry.

Applying mathlib's \lt{Ring.DirectLimit} to $\{A_k, \iota_k\}$ produces the
$\mathbb{C}$-algebra $R := \varinjlim_k A_k$, with a \texttt{StarRing}
structure descending from the entrywise conjugate transpose. Every element
of $R$ has a single-level representative: mathlib's
\lt{DirectLimit.exists\_eq\_mk} produces, from $x \in R$, a level $k$ and a
matrix $A \in A_k$ with $x = [\![(k,A)]\!]$. Because the embeddings are
isometric, the level norm descends to a well-defined norm on $R$ that makes
it a normed $*$-ring; the canonical map $\iota_k \colon A_k \to R$ is an
isometry, and the cocone identity holds up to \lt{Quotient.sound}.

The metric completion of $R$ under this norm is
\[
\Tinf \;:=\; \overline{\varinjlim_k A_k}^{\;\|\cdot\|_\op},
\]
built with \lt{UniformSpace.Completion}. The algebraic operations extend by
uniform continuity, and the $C^*$-identity passes to the completion by the
usual Cauchy-sequence argument, run against mathlib's completion API.

\begin{theorem}[UHF $C^*$-algebra registration]\label{thm:uhf-reg}
$\Tinf$ carries mathlib-native \texttt{Ring}, \texttt{StarRing},
\texttt{NormedRing}, and \texttt{CStarAlgebra} instances, with
$\varinjlim_k A_k$ dense in $\Tinf$.
\end{theorem}

By construction $\Tinf$ is a UHF algebra in the classical sense: the
norm-closure of an increasing union of full matrix algebras under unital
$*$-embeddings, of supernatural number $3^\infty$. The mathlib
\texttt{CStarAlgebra} instance is the algebra of record; no appeal is made
to the classical classification. Each level is itself simple, via
\lt{DivisionRing.isSimpleRing} on $\mathbb{C}$ composed with
\lt{IsSimpleRing.matrix}, and this finite-dimensional fact is the seed of
the simplicity argument in Section~\ref{sec:engine}. The base $3$ plays no
special role; every argument below works verbatim for any base $\ge 2$, and
the choice is inherited from the source project's ternary conventions.

%% =====================================================================
\section{The trace and its soft properties}\label{sec:trace}
%% =====================================================================

On level $k$ the normalized trace $\tau_n(M) := \Trace(M)/n$, with
$n = 3^k$, is linear, unital, and $1$-Lipschitz for the operator norm; the
Lipschitz bound uses the Hilbert--Schmidt estimate
$\|M\|_{\mathrm{HS}}^2 \le n\|M\|_\op^2$ and Cauchy--Schwarz. The traces are
compatible with the embeddings, $\tau_{3^{k+1}}(\iota_k A) = \tau_{3^k}(A)$,
by \lt{Matrix.trace\_kronecker}. The universal property of the direct limit
(\lt{DirectLimit.lift}) yields a pre-trace
$\tau_{\mathrm{pre}} \colon R \to \mathbb{C}$
(\lt{substrate\_pre\_trace}) that is additive, unital, and $1$-Lipschitz,
hence uniformly continuous, and
\[
\tuhf \;:=\; \lt{UniformSpace.Completion.extension}\,(\tau_{\mathrm{pre}})
\colon \Tinf \to \mathbb{C}
\]
is its unique continuous extension, agreeing with $\tau_{\mathrm{pre}}$ on
the dense image.

\begin{theorem}[The canonical tracial state]\label{thm:tracial}
$\tuhf \colon \Tinf \to \mathbb{C}$ is additive, unital, $1$-Lipschitz, and
uniformly continuous; $\mathbb{C}$-linear
(\lt{UHF\_trace\_smul}, bundled as \lt{UHF\_trace\_linearMap}); tracial,
$\tuhf(xy) = \tuhf(yx)$ (\lt{UHF\_trace\_mul\_comm}); Hermitian,
$\tuhf(x^*) = \overline{\tuhf(x)}$; and positive,
$0 \le \tuhf(x^*x)$ in the complex order
(\lt{UHF\_trace\_star\_mul\_self\_nonneg}, under mathlib's scoped
\texttt{ComplexOrder}).
\end{theorem}

Each property is a closed condition and follows the same three steps: a
mathlib matrix lemma at level $k$ (among \lt{Matrix.trace\_mul\_comm},
\lt{Matrix.trace\_conjTranspose}, and
\lt{Matrix.posSemidef\_conjTranspose\_mul\_self} with
\lt{Matrix.PosSemidef.trace\_nonneg}); reduction to a single-level
representative; and completion induction
(\lt{UniformSpace.Completion.induction\_on}, or its binary form for
traciality) over the closed set where the identity holds. Closedness of the
positivity set uses \lt{Complex.orderClosedTopology}.

Faithfulness stratifies over the three tiers of the construction. At the
matrix level, $\tau_n(M^*M) = 0 \Rightarrow M = 0$ is
\lt{Matrix.trace\_conjTranspose\_mul\_self\_eq\_zero\_iff}: the trace of
$M^*M$ is $\sum_{i,j}|M_{ij}|^2$. At the direct limit, reduce to a
single-level representative, apply the matrix case, and push the zero back
through the isometric embedding
(\lt{substrate\_pre\_trace\_star\_mul\_self\_faithful}). The completion
tier, $\tuhf(x^*x) = 0 \Rightarrow x = 0$ for $x \in \Tinf$, is the hard
case for the reason given in the introduction: it is not a closed
condition, and no density argument reaches it. Its proof occupies
Sections~\ref{sec:engine} and~\ref{sec:condexp}.

%% =====================================================================
\section{The simplicity engine and the false summit}\label{sec:engine}
%% =====================================================================

This section builds the algebraic machinery that first reduced
completion-tier faithfulness to a single hypothesis and then identified
that hypothesis as equivalent to the goal. All results are unconditional
kernel-verified theorems.

\subsection{Cauchy--Schwarz for the trace pairing}\label{sec:cs}

Define $\langle x, y \rangle_\tau := \tuhf(y^*x)$ on $\Tinf$
(\lt{traceInner}). Positivity (Theorem~\ref{thm:tracial}) makes this a
positive semidefinite sesquilinear form, and the discriminant argument goes
through against mathlib's \lt{discrim\_le\_zero}.

\begin{theorem}[\lt{UHF\_trace\_cauchySchwarz}]\label{thm:cs}
For all $x, y \in \Tinf$,
$\bigl\|\tuhf(y^*x)\bigr\|^2 \le \Re\tuhf(x^*x)\cdot \Re\tuhf(y^*y)$.
\end{theorem}

The proof expands $\tuhf\bigl((x - cy)^*(x - cy)\bigr) \ge 0$ along real
multiples $c = t\,\tuhf(y^*x)$, using linearity and the \texttt{StarModule}
structure, and applies the quadratic discriminant bound. One corollary is
left absorption of the trace-null cone
(\lt{UHF\_trace\_null\_left\_absorption}): if $\tuhf(x^*x) = 0$ then
$\tuhf\bigl((ax)^*(ax)\bigr) = 0$ for every $a$, since the left side is
$\langle a^*ax, x\rangle_\tau$, dominated through Theorem~\ref{thm:cs} by
$\Re\tuhf(x^*x) = 0$.

\subsection{The trace-null two-sided ideal}\label{sec:null}

Let $N := \{x \in \Tinf : \tuhf(x^*x) = 0\}$. Theorem~\ref{thm:cs} makes
$N$ additive; traciality makes it $*$-closed, since $\tuhf(xx^*) =
\tuhf(x^*x)$; left absorption and $*$-closure together give right
absorption; scalar closure is immediate.

\begin{theorem}[\lt{traceNullTwoSidedIdeal}]\label{thm:null}
$N$ is the carrier of a \texttt{TwoSidedIdeal} of $\Tinf$, and it is
proper, since $\tuhf(1) = 1$ gives $1 \notin N$
(\lt{traceNullTwoSidedIdeal\_ne\_top}).
\end{theorem}

This yields at once the conditional form of faithfulness, which was the
honest theorem available at that stage.

\begin{theorem}[\lt{UHF\_trace\_faithful\_of\_simple}]\label{thm:faithful-of-simple}
If every two-sided ideal of $\Tinf$ is $\{0\}$ or $\Tinf$, then $\tuhf$ is
faithful.
\end{theorem}

Indeed $N$ is a proper two-sided ideal, so simplicity forces $N = \{0\}$,
and $x \in N$ means $\tuhf(x^*x) = 0$, whence $x = 0$. The hypothesis is
genuine algebraic simplicity, classically Glimm's theorem.

\subsection{The Weyl averaging computation}\label{sec:weyl}

The engine of every simplicity proof for matrix towers is
finite-dimensional. Fix $n \ge 1$, let $\omega := e^{2\pi i/n}$ (a
primitive root, \lt{substrateRoot}), and let
\[
C := \diag(j \mapsto \omega^j), \qquad S_{ij} := [\,i = j+1\,]
\]
be the clock and shift matrices~\cite{weyl1927}, both unitary, with powers
$C^m = \diag(j \mapsto \omega^{jm})$ and $(S^m)_{ij} = [\,i = j+m\,]$. The
character sum $\sum_{a<n}(\omega^i\overline{\omega^j})^a = n\,[\,i=j\,]$
(\lt{substrate\_clock\_character\_sum}, via \lt{geom\_sum\_mul} and
\lt{IsPrimitiveRoot.pow\_inj}) drives a two-stage computation. Averaging
over the clock powers kills the off-diagonal,
$\tfrac{1}{n}\sum_a C^a x (C^a)^* = \diag(\diag x)$; averaging a diagonal
over the shift powers equalizes it,
$\tfrac{1}{n}\sum_b S^b D (S^b)^* = (\Trace D/n)\cdot 1$. Together:

\begin{theorem}[\lt{matrix\_unitary\_average\_eq\_trace\_smul\_one}]\label{thm:weyl}
For every $x \in M_n(\mathbb{C})$,
\[
\frac{1}{n^2}\sum_{a<n}\sum_{b<n}(C^aS^b)\,x\,(C^aS^b)^* \;=\;
\frac{\Trace x}{n}\cdot 1.
\]
\end{theorem}

A set closed under addition, arbitrary two-sided multiplication, and
$\mathbb{C}$-scaling absorbs each conjugate $uxu^*$, hence absorbs the
average: any such absorbing set containing $x$ contains
$(\Trace x/n)\cdot 1$.

\subsection{Simplicity below the completion}\label{sec:presimple}

Firing the engine at a nonzero element needs a nonzero trace, which a
nonzero ideal element need not have. The classical fix is the matrix-unit
trick: if $x \ne 0$, pick a nonzero entry $x_{pq}$; the matrix unit
$E_{q,p}$ satisfies $\Trace(E_{q,p}x) = x_{pq} \ne 0$
(\lt{trace\_single\_mul}), and $E_{q,p}x$ remains in the absorbing set.
Averaging and rescaling put $1$ into the set. With level simplicity this
gives matrix simplicity in ideal form
(\lt{matrix\_twoSidedIdeal\_bot\_or\_top}), and pulling a two-sided ideal
of $R$ back to each level (\lt{substrateIdealPullback}) lifts it to the
limit.

\begin{theorem}[\lt{directLimit\_twoSidedIdeal\_bot\_or\_top}]\label{thm:dl-simple}
Every two-sided ideal of $\varinjlim_k M_{3^k}(\mathbb{C})$ is $\{0\}$ or
the whole ring.
\end{theorem}

The same Weyl words push forward to unitaries $u_{a,b}$ in $\Tinf$, and the
completion-level averaging operator
$\lt{avgK}(x) := (3^k)^{-2}\sum_{a,b<3^k} u_{a,b}\,x\,u_{a,b}^*$ agrees with
$\tuhf(\cdot)\cdot 1$ on the image of level $k$, by Theorem~\ref{thm:weyl}.
Combining density, non-expansiveness of averaging, and the $1$-Lipschitz
property of $\tuhf$ gives, for every $x$ and $\varepsilon > 0$, a level $k$
with $\|\lt{avgK}(x) - \tuhf(x)\cdot 1\| < \varepsilon$ (\lt{avgK\_approx}).
Suppose a two-sided ideal $I$ contains an $x$ with $c := \tuhf(x) \ne 0$.
Averaging is by conjugation, so $\lt{avgK}(x) \in I$; taking
$\varepsilon = \|c\|$ gives $\|1 - c^{-1}\lt{avgK}(x)\| < 1$, so
$c^{-1}\lt{avgK}(x)$ is invertible by the Neumann series. An invertible
element of an ideal forces $1 \in I$: multiply it by its inverse, which
lies in the whole algebra, and the product stays in $I$ because $I$ absorbs
multiplication on both sides. No closure of $I$ enters: $\lt{avgK}(x)$ is a
finite average, a genuine element of $I$, so the argument applies to
algebraic, not merely closed, two-sided ideals. Thus $\Tinf$ is simple, and
$\tuhf$ is faithful, provided every nonzero ideal contains an element of
nonzero trace.

\subsection{The four-fold equivalence}\label{sec:tfae}

The natural next question is whether that trace-detection hypothesis can be
discharged directly. The formalized answer is no: the hypothesis is
logically equivalent to the goal it was meant to unlock.

\begin{theorem}[\lt{traceDetects\_tfae}]\label{thm:tfae}
For $\Tinf$ the following are equivalent:
\textup{(1)} every two-sided ideal $I \ne \{0\}$ contains an $x$ with
$\tuhf(x) \ne 0$;
\textup{(2)} $\lt{traceNullTwoSidedIdeal} = \{0\}$;
\textup{(3)} $\tuhf$ is faithful;
\textup{(4)} every two-sided ideal of $\Tinf$ is $\{0\}$ or $\Tinf$.
The cycle
$(4)\Rightarrow(3)\Rightarrow(2)\Rightarrow(1)\Rightarrow(4)$
is formalized with mathlib's \texttt{List.TFAE}.
\end{theorem}

This is the honest obstruction map. No hypothesis strictly weaker than
faithfulness closes the loop, and the classical hard input, Glimm's theorem
or something of equal strength, is genuinely required. The same development
also probed the spectral-projection route, which manufactures a nonzero
projection inside an ideal from a nonzero selfadjoint element through the
continuous functional calculus, and located exactly where it stalls in the
non-unital setting. The route that succeeded is the conditional
expectation, and it is different from both.

%% =====================================================================
\section{The level expectation and the faithfulness proof}\label{sec:condexp}
%% =====================================================================

The closing argument constructs the canonical projection onto each finite
level, proves two contraction estimates for it by elementary means, extends
it to the completion, and reads off faithfulness from the reconstruction
property. Simplicity then follows through Theorem~\ref{thm:tfae}. We write
$\tau_{3^k}$ for the normalized trace on $A_k = M_{3^k}(\mathbb{C})$.

\subsection{The partial trace}\label{sec:Ek}

Under $M_{3^{k+1}}(\mathbb{C}) \cong M_{3^k}(\mathbb{C}) \otimes
M_3(\mathbb{C})$, given by \lt{levelStepEquiv}, the single-step partial
trace $\mathrm{id}\otimes\tau_3 \colon A_{k+1} \to A_k$ is, entrywise,
\[
(\mathcal{E}B)_{ij} \;=\; \frac{1}{3}\sum_{t=0}^{2} B_{(i,t),(j,t)}
\qquad(\lt{partialTraceStep}),
\]
where $(i,t)$ abbreviates $\lt{levelStepEquiv}\,(i,t)$. Iterating gives
\lt{partialTraceDown}$\colon A_j\to A_k$ for $k \le j$.

These maps are compatible with the inductive system: for $k \le j \le j'$,
the iterated partial trace $A_{j'} \to A_k$ factors through $A_j$, and
$\mathcal{E}$ inverts a single embedding
($\mathcal{E}(A\otimes I_3) = A$, \lt{partialTraceStep\_substrateEmbedMatrix}).
Consequently the level maps agree on the image of every embedding, so they
descend, by the universal property of the direct limit, to a single
well-defined map
$E_k \colon \varinjlim_m A_m \to M_{3^k}(\mathbb{C})$ (\lt{condExp});
concretely, $E_k[\![(j,A)]\!] = \lt{partialTraceDown}\,(A)$ is independent
of the chosen single-level representative. This quotient-invariance is the
one nontrivial well-definedness check behind Theorem~\ref{thm:Ek}.

We call $E_k$ the \emph{level-$k$ expectation map}, and we use only the
retraction and trace-preservation properties recorded in
Theorem~\ref{thm:Ek} below. The full conditional-expectation API, meaning
positivity, the $A_k$-bimodule identity $E_k(axb) = a\,E_k(x)\,b$ for
$a,b \in A_k$, and idempotence as a norm-one projection, is not formalized.
The phrase ``conditional expectation'' is here shorthand for this
partial-trace retraction, not a claim that the classical axioms are
machine-checked.

\begin{theorem}[Properties of $E_k$]\label{thm:Ek}
$E_k$ is additive and $\mathbb{C}$-homogeneous, unital, star-preserving,
inverts the embedding on the level-$k$ corner
(\lt{condExp\_level\_k}, with \lt{condExp\_lower\_level} for levels $m \le
k$), and is trace-preserving:
$\tau_{3^k}(E_k x) = \tau_{\mathrm{pre}}(x)$
(\lt{condExp\_normalized\_trace}).
\end{theorem}

\subsection{The 2-norm contraction}\label{sec:l2}

The first analytic requirement is contraction of $E_k$ in the trace
$2$-seminorm:
\begin{equation}\label{eq:l2}
\Re\tau_{3^k}\bigl((E_k x)^*(E_k x)\bigr) \;\le\;
\Re\tau_{\mathrm{pre}}(x^*x)
\qquad(\lt{condExp\_l2\_contraction}).
\end{equation}
Classically this is Kadison--Schwarz for the completely positive map $E_k$,
composed with the trace. Mathlib has neither, and formalizing them would
have been a project in itself. The formalization targets the scalar
inequality directly, and it is elementary. Finite Cauchy--Schwarz over
three terms gives $\bigl|\sum_{t=0}^{2}a_t\bigr|^2 \le 3\sum_{t=0}^{2}|a_t|^2$
(\lt{normSq\_sum\_fin3\_le}), so each entry obeys
$|(\mathcal{E}B)_{ij}|^2 \le \tfrac{1}{3}\sum_t|B_{(i,t),(j,t)}|^2$. Summing
over $i,j$, the right side is exactly the diagonal-tail sub-sum ($s = t$) of
\[
\tau_{3^{k+1}}(B^*B) \;=\; \frac{1}{3^{k+1}}\sum_{i,s,j,t}\bigl|B_{(j,t),(i,s)}\bigr|^2,
\]
every term of which is nonnegative. Mathlib's \lt{Finset.single\_le\_sum}
closes the single-step bound
$\tau_{3^k}\bigl((\mathcal{E}B)^*(\mathcal{E}B)\bigr) \le \tau_{3^{k+1}}(B^*B)$
(\lt{partialTraceStep\_normSq\_trace\_le}); iterating and descending to the
limit gives~\eqref{eq:l2}. No positive-semidefinite operator theory and no
complete positivity enter. For trace-compatible partial traces the scalar
$2$-norm contraction is a three-line fact, and the Kadison--Schwarz detour
of the textbook proof is avoidable.

\subsection{The operator-norm contraction}\label{sec:opnorm}

Extending $E_k$ to the completion requires the operator-norm bound
$\|E_k x\| \le \|x\|$. Classically this is the statement that a unital
completely positive map has norm one. The formalized proof uses an isometry
decomposition. For $t \in \{0,1,2\}$ let
$W_t \colon \mathbb{C}^{3^k} \to \mathbb{C}^{3^{k+1}}$ be the coordinate
injection $v \mapsto v \otimes e_t$, viewed as a
$3^{k+1}\times 3^k$ matrix (\lt{Winj}). Then
\begin{equation}\label{eq:stinespring}
\mathcal{E}B \;=\; \frac{1}{3}\sum_{t=0}^{2} W_t^*\, B\, W_t
\qquad(\lt{partialTraceStep\_eq\_smul\_sum}),
\end{equation}
and each $W_t$ satisfies $W_t^*W_t = I$. The $C^*$-identity for the
$L^2$-operator norm gives $\|W_t\|^2 = \|W_t^*W_t\| = \|I\| \le 1$
(\lt{Winj\_opNorm\_sq\_le\_one}), so submultiplicativity and
$\|W_t^*\| = \|W_t\|$ yield $\|W_t^*BW_t\| \le \|B\|$, and averaging over
the three slots gives $\|\mathcal{E}B\| \le \|B\|$
(\lt{partialTraceStep\_opNorm\_le}). Iterating gives $\|E_kx\| \le \|x\|$
(\lt{condExp\_opNorm\_le}), so $E_k$ is $1$-Lipschitz
(\lt{condExp\_lipschitz}), hence uniformly continuous. The only inputs are
four $L^2$-operator-norm lemmas and the norm bound on the identity operator.
Decomposition~\eqref{eq:stinespring} is the concrete Stinespring form of a
partial trace~\cite{stinespring1955}.

\subsection{Extension and the closing argument}\label{sec:summit}

Uniform continuity extends $E_k$ to the completion in two forms: the
matrix-valued \lt{condExpFin}~$k\colon \Tinf \to M_{3^k}(\mathbb{C})$, and
its recoordinatization \lt{condExpCompletion}~$k\colon \Tinf \to \Tinf$
given by $x \mapsto \iota_k(\lt{condExpFin}\,k\,x)$. Both are $1$-Lipschitz
and agree with $E_k$ on the dense image. Three statements complete the
proof.

\begin{lemma}[Reconstruction, \lt{condExpCompletion\_tendsto}]\label{lem:tendsto}
For every $x \in \Tinf$, $\lt{condExpCompletion}\,k\,x \to x$ as
$k \to \infty$.
\end{lemma}

The proof is an $\varepsilon/2$ density lift over the uniformly Lipschitz
family: approximate $x$ within $\varepsilon/2$ by a dense-image element with
representative at level $m$; for $k \ge m$ the map fixes that element
(\lt{condExpCompletion\_fixes\_coe\_lower}), and the Lipschitz bound
controls the difference at $x$.

\begin{lemma}[Lifted contraction, \lt{condExpCompletion\_l2\_contraction}]\label{lem:l2-lift}
For every $k$ and $x \in \Tinf$,
$\Re\tuhf\bigl((\lt{condExpCompletion}\,k\,x)^*(\lt{condExpCompletion}\,k\,x)\bigr)
\le \Re\tuhf(x^*x)$.
\end{lemma}

This is an inequality between two continuous real-valued functions of $x$.
It holds on the dense image by~\eqref{eq:l2} and the trace bridge
$\tuhf\bigl((E_kx)^*(E_kx)\bigr) =
\tau_{3^k}\bigl((\lt{condExpFin}\,k\,x)^*(\lt{condExpFin}\,k\,x)\bigr)$,
and completion induction on the closed set $\{x : f(x) \le g(x)\}$
(\lt{isClosed\_le}) transfers it to all of $\Tinf$.

\begin{lemma}[Matrix faithfulness, real-part form,
\lt{matrix\_faithful\_of\_re\_nonpos}]\label{lem:re}
If $\Re\tau_n(M^*M) \le 0$ then $M = 0$.
\end{lemma}

Positivity pins $\tau_n(M^*M)$ to the nonnegative reals; a nonpositive real
part then forces it to $0$, and matrix-tier faithfulness concludes.

\begin{theorem}[Faithfulness, \lt{UHF\_trace\_faithful}]\label{thm:faithful}
For every $x \in \Tinf$, if $\tuhf(x^*x) = 0$ then $x = 0$. Kernel axioms
$[\texttt{propext}, \texttt{Classical.choice}, \texttt{Quot.sound}]$; zero
project axioms.
\end{theorem}

\begin{proof}
Suppose $\tuhf(x^*x) = 0$. By Lemma~\ref{lem:l2-lift} and the trace bridge,
$\Re\tau_{3^k}\bigl((\lt{condExpFin}\,k\,x)^*(\lt{condExpFin}\,k\,x)\bigr)
\le 0$ for every $k$; by Lemma~\ref{lem:re},
$\lt{condExpFin}\,k\,x = 0$, hence $\lt{condExpCompletion}\,k\,x = 0$, for
every $k$. By Lemma~\ref{lem:tendsto} the constant-zero sequence
$\lt{condExpCompletion}\,k\,x$ converges to $x$, so uniqueness of limits
gives $x = 0$.
\end{proof}

\begin{corollary}[Simplicity, \lt{substrate\_completion\_simple\_unconditional}]\label{cor:simple}
Every two-sided ideal of $\Tinf$ is $\{0\}$ or $\Tinf$.
\end{corollary}

\begin{proof}
Theorem~\ref{thm:faithful} is item (3) of Theorem~\ref{thm:tfae}, and item
(4) follows.
\end{proof}

\begin{remark}
The order of deduction is the reverse of Glimm's: he proves simplicity
first and derives faithfulness of the trace, whereas here faithfulness is
proved directly through the level expectations and simplicity follows. Both
directions of the equivalence are formalized, so Glimm's implication,
Theorem~\ref{thm:faithful-of-simple}, is machine-checked too.
Corollary~\ref{cor:simple} is the stronger algebraic statement, covering
all two-sided ideals rather than only closed ones, which holds for every
UHF algebra~\cite{davidson1996}.
\end{remark}

\subsection{Uniqueness of the trace}\label{sec:unique}

The same level-expectation machinery gives the third leg of the
identification: $\tuhf$ is the only trace. We package the hypotheses as a
predicate \lt{IsTracialState}~$\varphi$ requiring $\varphi$ to be
continuous, additive, $\mathbb{C}$-homogeneous, unital, and tracial;
positivity is deliberately omitted. (The Lean name retains ``State'' for
continuity with the source; because positivity is dropped, the predicate is
really that of a continuous unital tracial \emph{functional}, and the
theorem is correspondingly stronger than uniqueness among states.)

\begin{theorem}[Uniqueness, \lt{substrate\_UHF\_trace\_unique}]\label{thm:unique}
For every $\varphi$ with \lt{IsTracialState}~$\varphi$ and every
$x \in \Tinf$, $\varphi(x) = \tuhf(x)$. Kernel axioms
$[\texttt{propext}, \texttt{Classical.choice}, \texttt{Quot.sound}]$.
\end{theorem}

The proof is the standard argument in three steps. At a finite level, any
additive, $\mathbb{C}$-homogeneous, unital, tracial functional on
$M_n(\mathbb{C})$ equals the normalized trace
(\lt{matrix\_tracial\_state\_unique}): the off-diagonal matrix units are
commutators, $E_{ij} = [E_{ik}, E_{kj}]$, hence annihilated, and the
diagonal units are cyclic-conjugate, hence equal, so $\varphi$ is $\varphi(1)/n$
times the trace, and unitality fixes the constant. Mathlib does not carry
this lemma, so it is proved from \lt{Matrix.single} identities and
\lt{matrix\_eq\_sum\_single}. For $x = \iota_k(A)$, the functional
$\varphi\circ\iota_k$ is of this kind, so $\varphi(x) = \tau_{3^k}(A) =
\tuhf(x)$. Finally $\varphi$ and $\tuhf$ are continuous and agree on the
dense image, so they agree everywhere, by \lt{isClosed\_eq}.

Theorem~\ref{thm:faithful}, Corollary~\ref{cor:simple}, and
Theorem~\ref{thm:unique} together record, in the capstone
\lt{r113\_substrate\_UHF\_factor\_capstone}, that $\Tinf$ is a faithful,
simple, uniquely-traced UHF $C^*$-algebra: the Glimm $3^\infty$ factor.

\subsection{Mathlib infrastructure and gaps}\label{sec:formalization}

The construction is mathlib-native. The load-bearing APIs are
\lt{Ring.DirectLimit} with \lt{DirectLimit.lift} and
\lt{DirectLimit.exists\_eq\_mk}; \lt{UniformSpace.Completion} with its
\lt{extension}, \lt{induction\_on}, and isometry lemmas; the scoped
$L^2$-operator-norm $C^*$-structure on matrices; the matrix trace and
positivity lemmas; \lt{IsSimpleRing.matrix} and \lt{TwoSidedIdeal}; the
roots-of-unity API for the Weyl computation; and the scoped
\texttt{ComplexOrder} with \lt{discrim\_le\_zero}.

Three absences shaped the proofs and are natural next targets for mathlib.
There is no complete-positivity theory: no Kadison--Schwarz inequality, no
partial trace, no ``unital CP map has norm one.'' This forced, and in the
end rewarded, the elementary routes of Sections~\ref{sec:l2}
and~\ref{sec:opnorm}, which are shorter than the classical arguments and
may be of expository interest on their own. There are no spectral
projections: the Borel functional calculus is absent, and the continuous
functional calculus produces projections only from clopen spectral sets;
the probe of Section~\ref{sec:tfae} built exactly that and identified the
missing hypothesis. And a \lt{UniformSpace.Completion} does not
automatically carry the order structure of a $C^*$-algebra, so positivity
of $\tuhf$ was proved scalar-side in the \texttt{ComplexOrder} on
$\mathbb{C}$ rather than operator-side. Three standalone files written to
mathlib contribution standards were produced during the spectral-projection
probe, on topological closures of two-sided ideals, on the non-unital
continuous functional calculus mapping into closed ideals, and on spectral
projections from clopen quasispectral sets, and are offered upstream.

%% =====================================================================
\section{Reproducibility}\label{sec:repro}
%% =====================================================================

\paragraph{Pinned toolchain.}
Lean toolchain \texttt{leanprover/lean4:v4.24.0-rc1}; mathlib commit
\lt{eed770a434957369c6262aa3fb1d6426419016d4}; verified HEAD
\texttt{91cc1a01} (2026-07-23), with build and \texttt{\#print axioms} run
the same day. The source is public at
\href{https://github.com/FractalDevTeam/Principia-Fractalis}{\lt{github.com/FractalDevTeam/Principia-Fractalis}};
the commit above is a stable reference point in that history.

\paragraph{Build.}
From the \texttt{PF\_Lean4\_Code} directory, \lt{elan install \$(cat
lean-toolchain)} followed by \lt{lake build PF} (4{,}472 jobs) and
\lt{lake build} (8{,}930 jobs; the second pass re-elaborates every theorem
through the production Lean~4 kernel under a distinct package boundary).

\paragraph{Axiom audit.}
The audit blocks at the ends of the source files report, for each headline
theorem:

\begin{small}
\begin{verbatim}
'PrincipiaTractalis.SubstrateCompletionFaithful.UHF_trace_faithful'
  depends on axioms: [propext, Classical.choice, Quot.sound]

'PrincipiaTractalis.SubstrateCompletionFaithful.
  substrate_completion_simple_unconditional'
  depends on axioms: [propext, Classical.choice, Quot.sound]

'PrincipiaTractalis.SubstrateTraceUniqueness.substrate_UHF_trace_unique'
  depends on axioms: [propext, Classical.choice, Quot.sound]
\end{verbatim}
\end{small}

These are the three standard Lean~4 kernel axioms: propositional
extensionality, choice, and quotient soundness. There are no project-level
axioms and no \texttt{sorry} on any code path leading to the theorems of
this paper. The development is classical, not constructive: choice enters
through mathlib's extension of uniformly continuous maps to a completion
(\lt{UniformSpace.Completion.extension}) and the universal properties of
the direct limit, both used pervasively here. We make no claim of
computable or choice-free content.

\begin{table}[ht]
\centering
\small
\begin{tabular}{@{}lll@{}}
\toprule
Content & Section & Lean file (under \texttt{PF/}) \\
\midrule
Tower, limit, completion, instances & \S\ref{sec:tower} & \lt{Substrate*} chain \\
Trace and soft properties & \S\ref{sec:trace} & \lt{SubstrateUHFTraceIs*.lean} \\
Cauchy--Schwarz for the trace & \S\ref{sec:cs} & \lt{SubstrateUHFTraceCauchySchwarz.lean} \\
Trace-null ideal; faithful-of-simple & \S\ref{sec:null} & \lt{SubstrateUHFTraceNullIdeal.lean} \\
Weyl averaging & \S\ref{sec:weyl} & \lt{SubstrateMatrixUnitaryAveraging.lean} \\
Direct-limit simplicity & \S\ref{sec:presimple} & \lt{SubstrateDirectLimitSimplicity.lean} \\
Four-fold equivalence & \S\ref{sec:tfae} & \lt{SubstrateTraceDetectionAttempt.lean} \\
Level expectation $E_k$ & \S\ref{sec:Ek} & \lt{SubstrateConditionalExpectation.lean} \\
$2$-norm contraction & \S\ref{sec:l2} & \lt{SubstrateCondExpContraction.lean} \\
Operator-norm contraction & \S\ref{sec:opnorm} & \lt{SubstrateCondExpOpNorm.lean} \\
Extension; faithfulness; simplicity & \S\ref{sec:summit} & \lt{SubstrateCompletionFaithful.lean} \\
Uniqueness; factor capstone & \S\ref{sec:unique} & \lt{SubstrateTraceUniqueness.lean} \\
\bottomrule
\end{tabular}
\medskip
\caption{Section-to-source map.}\label{tab:files}
\end{table}

%% =====================================================================
\section{Related work}\label{sec:related}
%% =====================================================================

Mathlib's $C^*$-algebra hierarchy, continuous functional calculus, matrix
positivity infrastructure, direct limits, and completions, all community
developments with functional-analysis foundations described
in~\cite{dupuis-lewis-macbeth2022}, are the platform this work builds on.
To the best of our knowledge, after searching mathlib, the Coq
\texttt{opam} package archive, and the Isabelle Archive of Formal Proofs
for UHF, AF, Glimm, simple-$C^*$, and tracial-state content, and after
consulting the operator-algebra formalization literature, no prior
formalization in Lean, Coq, or Isabelle contains a concrete
infinite-dimensional simple $C^*$-algebra, a canonical faithful tracial
state on one, or any version of Glimm's simplicity
theorem~\cite{glimm1960}. We have not performed an exhaustive survey of
every proof-assistant library and would welcome pointers to prior art.

On the classical side, everything proved here is textbook
material~\cite{davidson1996, blackadar2006, takesaki1979}: Glimm's
classification and simplicity~\cite{glimm1960}, Powers' work on the
associated factors~\cite{powers1967}, Bratteli's AF
framework~\cite{bratteli1972}, and the conditional-expectation theory of
Tomiyama~\cite{tomiyama1957}, descending from Kadison~\cite{kadison1952}
and Stinespring~\cite{stinespring1955}. The contribution is not new
mathematics about $M_{3^\infty}$; it is the first machine-checked path to
these facts, together with the elementary replacement proofs of
Sections~\ref{sec:l2} and~\ref{sec:opnorm} that the setting demanded.

Several extensions are natural. The three contributed files can be
upstreamed, and the tower can be de-specialized from base $3$ to an
arbitrary supernatural number, since most proofs are already
base-generic in structure. A general complete-positivity API, with partial
traces and Kadison--Schwarz, now has a worked concrete example to test
against. The GNS representation of $\tuhf$ and the trace on the double
commutant would connect this construction to the hyperfinite II$_1$
factor~\cite{murray-vonneumann1936, powers1967}, and Glimm's classification
by supernatural number is a further target for which the clopen spectral
projection API is a first ingredient.

\section*{Acknowledgments}
The Lean~4 and mathlib community built the $C^*$-algebra, direct-limit,
completion, and functional-calculus infrastructure on which every theorem
here stands.

\begin{thebibliography}{99}

\bibitem{blackadar2006}
B.~Blackadar.
\newblock \emph{Operator Algebras: Theory of $C^*$-Algebras and von Neumann Algebras}.
\newblock Encyclopaedia of Mathematical Sciences 122, Springer, 2006.

\bibitem{bratteli1972}
O.~Bratteli.
\newblock Inductive limits of finite dimensional $C^*$-algebras.
\newblock \emph{Trans. Amer. Math. Soc.}, 171:195--234, 1972.

\bibitem{davidson1996}
K.~R.~Davidson.
\newblock \emph{$C^*$-Algebras by Example}.
\newblock Fields Institute Monographs 6, American Mathematical Society, 1996.

\bibitem{dupuis-lewis-macbeth2022}
F.~Dupuis, R.~Y.~Lewis, and H.~Macbeth.
\newblock Formalized functional analysis with semilinear maps.
\newblock In \emph{13th International Conference on Interactive Theorem Proving (ITP 2022)}, LIPIcs 237, 2022.

\bibitem{glimm1960}
J.~Glimm.
\newblock On a certain class of operator algebras.
\newblock \emph{Trans. Amer. Math. Soc.}, 95:318--340, 1960.

\bibitem{kadison1952}
R.~V.~Kadison.
\newblock A generalized Schwarz inequality and algebraic invariants for operator algebras.
\newblock \emph{Ann. of Math.~(2)}, 56:494--503, 1952.

\bibitem{mathlib2020}
The mathlib Community.
\newblock The Lean Mathematical Library.
\newblock In \emph{Proc.\ CPP 2020}, pages 367--381, 2020.

\bibitem{moura-ullrich2021}
L.~de Moura and S.~Ullrich.
\newblock The Lean~4 Theorem Prover and Programming Language.
\newblock In \emph{CADE 28}, 2021.

\bibitem{murray-vonneumann1936}
F.~J.~Murray and J.~von Neumann.
\newblock On rings of operators.
\newblock \emph{Ann. of Math.~(2)}, 37(1):116--229, 1936.

\bibitem{powers1967}
R.~T.~Powers.
\newblock Representations of uniformly hyperfinite algebras and their associated von Neumann rings.
\newblock \emph{Ann. of Math.~(2)}, 86:138--171, 1967.

\bibitem{stinespring1955}
W.~F.~Stinespring.
\newblock Positive functions on $C^*$-algebras.
\newblock \emph{Proc. Amer. Math. Soc.}, 6:211--216, 1955.

\bibitem{takesaki1979}
M.~Takesaki.
\newblock \emph{Theory of Operator Algebras I}.
\newblock Springer, 1979.

\bibitem{tomiyama1957}
J.~Tomiyama.
\newblock On the projection of norm one in $W^*$-algebras.
\newblock \emph{Proc. Japan Acad.}, 33:608--612, 1957.

\bibitem{weyl1927}
H.~Weyl.
\newblock Quantenmechanik und Gruppentheorie.
\newblock \emph{Zeitschrift f\"ur Physik}, 46:1--46, 1927.

\end{thebibliography}

\end{document}
